\begin{onehalfspace}

\section[Ottimizzazione topologica della QFT]{Ottimizzazione della trasformata di Fourier consapevole della topologia}
\label{sec:qft_topologia}

La QFT approssimata elimina alcune rotazioni per ridurre il costo del circuito. La campagna esplorativa misura se il rumore evitato compensi la perdita di precisione. I test distinguono Shor dalla sola stima di fase, perché i due circuiti hanno costi molto diversi.

\paragraph{Controllo su $N=15$ ed eliminazione degli SWAP.}
Il pilota su $N=15$ sceglie $k_{\mathrm{QPE}}=1$ sui dati di selezione. Sui dati indipendenti di verifica, la variante ottiene $0.6279$ contro $0.4792$ della QFT piena. La differenza è $+0.1487$, con IC Newcombe 95\% $[+0.1273;+0.1699]$; le porte ECR scendono da 362 a 279. È però un caso favorevole, con fasi esatte in due bit, e non dimostra lo stesso beneficio per ordini più lunghi. Togliere gli SWAP che invertono i bit non dà invece un risparmio: conteggi compilati e risultati restano identici, perché il transpiler assorbe già la permutazione nel layout. Griglia e controllo della variante corretta sono nella Sezione~\ref{sec:app_aqft} dell'Appendice~\ref{app:fase2}.

Nell'aritmetica di Beauregard, gran parte del costo viene dalle QFT interne agli addizionatori modulari. Per $N=21$ si contano 11\,436 porte \texttt{cp} e 3\,406 \texttt{cx}; la QFT finale richiede invece solo 28 rotazioni, su circa 47\,000 porte complessive. Eliminare le rotazioni più piccole negli addizionatori porta le porte a due qubit da 49\,661 a 23\,178. Agire sulla sola trasformata finale non cambia l'ordine di grandezza. Il primo controllo identifica quindi la parte del circuito su cui conviene intervenire.

Ridurre le porte comporta una perdita di precisione. Su $N=21$, il successo ideale scende da $0{,}4667$ a $0{,}4013$, cioè di circa sei punti e mezzo. In presenza di rumore, la variante troncata può recuperare questa perdita solo se il risparmio di porte evita abbastanza errori. Il bilancio dipende quindi dalla qualità delle porte e va verificato.

Un modello semplificato combina il successo ideale con quello di una distribuzione uniforme. In questo modello, il vantaggio massimo stimato della variante troncata è circa \textbf{tre punti percentuali}. Richiederebbe errori a due qubit circa \textbf{duecento volte inferiori} a quelli dei dati di calibrazione considerati. È una stima del modello, non un vantaggio misurato né un limite universale. Nei test rumorosi non emerge un beneficio statisticamente significativo del troncamento. Al rumore nominale, le configurazioni provate restano vicine al livello uniforme; nel punto in cui compare un segnale, è la QFT \emph{piena} ad avere il risultato migliore. Il circuito di $N=21$ dura circa centocinquanta volte il tempo di coerenza tipico. Questa profondità limita la possibilità di distinguere le varianti nel regime considerato.

Per capire se il limite dipenda dal costo di Shor, si ripete il confronto sulla sola \textbf{stima di fase}. Con un unitario di fase, il circuito richiede decine di porte a due qubit, anziché decine di migliaia. I nove confronti usano tre dimensioni di registro e fasi sia esattamente rappresentabili sia non rappresentabili in $n$ bit. La variante troncata ha un risultato migliore in \textbf{tutti e nove}; in sei, l'intervallo di confidenza è interamente positivo. Su dieci qubit il vantaggio raggiunge $+0{,}126$, con IC Newcombe 95\% $[+0{,}098;\ +0{,}153]$.

Nei test sulla stima di fase emergono tre regolarità. Il grado migliore di troncamento resta \textbf{due o tre} alle taglie provate, mentre la QFT piena usa $n-1$. Il risparmio cresce così da venticinque porte a due qubit su sei qubit a settantotto su dieci. Cresce anche il vantaggio, da circa due a tredici punti percentuali. Infine, le fasi esattamente rappresentabili beneficiano più delle altre. Questo aiuta a interpretare la differenza fra le istanze di Shor: per $N=15$, l'ordine $r=4$ dà fasi $0$, $1/4$, $1/2$, $3/4$, esatte in due bit. Per $N=21$, l'ordine $r=6$ dà invece fasi $j/6$, per le quali le rotazioni fini conservano informazione utile.

Un limite teorico aiuta a interpretare il grado piccolo osservato, senza estenderlo a registri arbitrariamente grandi. Akoramurthy e Surendiran hanno derivato un limite sulla distanza di variazione totale fra la distribuzione QPE esatta e quella ottenuta con la trasformata troncata~\cite{akoramurthy2026}:
%
\begin{equation}
\mathrm{TV}\!\left(P_\phi,\, P_\phi^{d}\right) \;\le\; \frac{\pi\,(m-d)}{2^{d}},
\label{eq:bound_troncamento}
\end{equation}
%
dove $m$ è la dimensione del registro e $d$ il grado di troncamento. Il numeratore cresce linearmente in $m$ e il denominatore esponenzialmente in $d$. Per mantenere fisso il limite, $d$ deve quindi crescere circa logaritmicamente con il registro. Il grado due o tre osservato nei test non basta a dimostrare che un grado costante sia sempre sufficiente. Gli autori riportano una riduzione delle porte del $31{,}3$--$43{,}7\%$ e un costo che passa da $O(m^2)$ a $O(m \log m)$ per $d = O(\log m)$. Studiano però la stima di fase, senza la fattorizzazione. Il confronto con Shor e la scelta del grado su dati separati da quelli di verifica appartengono alla presente campagna. Su Shor, i circa tre punti percentuali restano una stima del modello ausiliario.

Conteggi di porte e calibrazione permettono di proporre gradi candidati
prima dell'esecuzione. In questa campagna, però, il grado viene scelto
empiricamente sui batch di selezione e verificato su holdout disgiunti:
non è stata validata una scelta ottima dai soli conteggi. Il contributo
è quindi un criterio di progetto e un protocollo di verifica per la
trasformata, da controllare quando cambia l'algoritmo o il dispositivo.

Un'estensione consiste nel combinare tre riduzioni. La prima, misurata qui, approssima la QFT per usare meno \emph{porte}. La seconda approssima l'aritmetica modulare per usare meno \emph{qubit di lavoro}~\cite{chevignard2025} e contribuisce alle recenti stime di risorse~\cite{gidney2025}. La Tabella~\ref{tab:traiettoria_risorse} ne riporta l'evoluzione e le assunzioni. La terza divide la stima di fase in blocchi indipendenti per ridurre i \emph{qubit di conteggio}~\cite{shukla2026}.

I tre risparmi non si sommano automaticamente. Ridurre i qubit di lavoro può richiedere più porte: per RSA-2048, una proposta passa da 6144 a 1730 qubit logici, ma porta i Toffoli da $2^{31{,}3}$ a $2^{40{,}9}$~\cite{chevignard2025}. Ottimizzazioni aritmetiche successive recuperano parte di questo costo~\cite{gidney2025}. Combinare la riduzione dei qubit con il troncamento della QFT richiede quindi di misurare il bilancio finale delle porte.

La riduzione del registro di conteggio agisce invece sulla stima di fase e lascia invariato il registro di lavoro~\cite{shukla2026}. La proposta passa da $2n+1$ a tre o quattro qubit di conteggio. Blocchi indipendenti stimano più bit di fase, poi un procedimento dal bit meno significativo al più significativo risolve le ambiguità~\cite{shuklaqpe2026}. Per verificare se questa riduzione sia compatibile con le altre, occorre un esperimento che le combini e misuri anche la profondità. Le singole proposte sono quantificate; il loro beneficio congiunto resta da dimostrare.

La QFT compare anche nella stima di fase, nell'amplitude estimation, nel conteggio quantistico, nella risoluzione di sistemi lineari e nel problema del sottogruppo nascosto, di cui Shor è un caso particolare. Si usa inoltre nelle simulazioni che alternano rappresentazioni in posizione e in momento. Questi algoritmi impiegano le stesse rotazioni controllate di piccolo angolo. Il protocollo seguito qui può suggerire verifiche analoghe: proporre gradi di troncamento dai conteggi, sceglierli su dati di selezione e valutarli su dati indipendenti. Queste verifiche restano da eseguire. Il contributo è quindi un metodo di valutazione, con risultati limitati ai circuiti provati.

Il pilota incompleto su $N=21$, i tre livelli di sensibilità al rumore,
la distinzione fra configurazione scelta e variante troncata, e i nove
contrasti QPE sono documentati nelle Tabelle~\ref{tab:aqft_n21_sensibilita}
e~\ref{tab:aqft_qpe_completa}. Sono riportati anche gli esiti negativi
e i contrasti i cui intervalli includono lo zero.

La verifica usa un layout per dimensione, una sola calibrazione e tre fasi bersaglio. Le varianti condividono batch e semi, quindi il confronto interno usa gli stessi campioni. Il risultato non dimostra però che il beneficio si mantenga su altri dispositivi o topologie.

Il precedente più vicino è la compilazione \emph{noise-adaptive}, che ordina o vincola il
collocamento del circuito usando i dati di calibrazione~\cite{murali2019, tannu2019}. La
differenza rispetto a quella linea, e il motivo per cui resta spazio, è che tali metodi
ottimizzano la fedeltà del circuito, mentre per un algoritmo come Shor conta soltanto che
il post-processing restituisca i fattori. I due obiettivi non coincidono, e la campagna
sulla compilazione riportata in Appendice~\ref{app:fase2} misura quanto non coincidano.

% -----------------------------------------------

\end{onehalfspace}
